\documentclass[10pt]{article}
\usepackage[T1]{fontenc}
\usepackage[utf8]{inputenc}
\usepackage{lmodern}
\renewcommand{\familydefault}{\sfdefault}
\usepackage[letterpaper,margin=0.60in]{geometry}
\usepackage{enumitem}
\usepackage{tabularx}
\usepackage{titlesec}
\usepackage{xcolor}
\usepackage{graphicx}
\usepackage{microtype}
\usepackage{fancyhdr}
\usepackage{hyperref}
\input{glyphtounicode}
\pdfgentounicode=1

\definecolor{accent}{HTML}{274690}
\definecolor{muted}{HTML}{51565E}
\definecolor{rulegray}{HTML}{B8BCC2}
\hypersetup{colorlinks=true,urlcolor=accent,linkcolor=accent,citecolor=accent,pdfborder={0 0 0}}
\setlength{\parindent}{0pt}
\setlength{\tabcolsep}{0pt}
\raggedright
\urlstyle{same}
\emergencystretch=1.5em

\titleformat{\section}{\large\bfseries}{}{0pt}{}
\titlespacing*{\section}{0pt}{8pt}{4pt}

% Optional portrait. Upload photo.jpg and set \showphototrue in the main file.
\newif\ifshowphoto
\showphotofalse
\newcommand{\photofile}{photo.jpg}

\newcommand{\contactline}{%
  \href{mailto:tharunv.pvit@gmail.com}{tharunv.pvit@gmail.com} $\vert$
  240-791-1188 $\vert$
  \href{https://tvpian.github.io/}{Portfolio} $\vert$
  \href{https://www.linkedin.com/in/tvpian/}{LinkedIn} $\vert$
  \href{https://github.com/tvpian}{GitHub} $\vert$
  \href{https://scholar.google.com/citations?user=Vre9wQQAAAAJ\&hl=en}{Scholar}%
}

\newcommand{\cvheaderplain}[1]{%
  \begin{center}
    {\LARGE\bfseries Tharun V. Puthanveettil}\\[3pt]
    {\small\contactline}\\[5pt]
    {\normalsize\bfseries #1}
  \end{center}
}

\newcommand{\cvheader}[1]{%
  \ifshowphoto
    \IfFileExists{\photofile}{%
      \begin{tabularx}{\textwidth}{@{}X r@{}}
        \begin{minipage}[c]{0.82\textwidth}
          {\LARGE\bfseries Tharun V. Puthanveettil}\\[3pt]
          {\small\contactline}\\[5pt]
          {\normalsize\bfseries #1}
        \end{minipage} &
        \raisebox{-0.5\height}{\includegraphics[width=0.78in,height=0.78in,keepaspectratio]{\photofile}}
      \end{tabularx}
    }{\cvheaderplain{#1}}
  \else
    \cvheaderplain{#1}
  \fi
}

\newcommand{\roleentry}[4]{%
  \begin{tabularx}{\textwidth}{@{}X r@{}}
    \textbf{#1} & \textbf{#2}\\
    \textit{#3} & \textit{#4}
  \end{tabularx}\vspace{-2pt}
}

\newenvironment{tightitems}{%
  \begin{itemize}[leftmargin=15pt,itemsep=2.1pt,topsep=3pt,parsep=0pt,partopsep=0pt]
}{\end{itemize}\vspace{1pt}}

\newenvironment{compactitems}{%
  \begin{itemize}[leftmargin=15pt,itemsep=1.3pt,topsep=2pt,parsep=0pt,partopsep=0pt]
}{\end{itemize}\vspace{0pt}}

\newcommand{\projectentry}[3]{%
  \begin{tabularx}{\textwidth}{@{}X r@{}}
    \textbf{#1} $\vert$ \textit{#2} & #3
  \end{tabularx}\vspace{-2pt}
}

\newcommand{\researchentry}[3]{%
  \begin{tabularx}{\textwidth}{@{}X r@{}}
    \textbf{#1} & \status{#2}\\[-1pt]
    \multicolumn{2}{@{}l@{}}{\textit{#3}}
  \end{tabularx}\vspace{-2pt}
}

\newcommand{\evidence}[2]{\nobreak\hspace{0.15em}{\footnotesize\href{#1}{[\MakeUppercase{#2}]}}}
\newcommand{\status}[1]{\textcolor{muted}{\footnotesize\MakeUppercase{#1}}}

\newcommand{\pageheader}[1]{%
  \fancyhead[L]{\small Tharun V. Puthanveettil}%
  \fancyhead[R]{\small #1}%
  \renewcommand{\headrulewidth}{0.3pt}%
}

\hypersetup{pdftitle={Tharun V. Puthanveettil – Research CV},pdfauthor={Tharun V. Puthanveettil}}
\pagestyle{fancy}\fancyhf{}\pageheader{Research CV}\cfoot{\small\thepage}
\newcommand{\thread}[3]{\par\vspace{4pt}\textbf{#1}\hfill{\scriptsize #2}\par\textit{#3}\par}
\begin{document}
\cvheader{Human-Inspired Robot Learning | Multimodal Learning | Human–Robot Interaction}
\fontsize{10}{12}\selectfont
\section{Research Focus}
I want to build robots that understand the people they work with and keep learning from them over time. I take inspiration from how people learn: they notice what they cannot yet do, learn it from the right person through demonstration, correction, or instruction, and keep what matters. I am interested in robots that grow the same way, adding small specialized skills on top of an understanding that keeps developing, and in what this requires of robot learning architectures. So far I have worked on multimodal understanding of people, learning from human intervention, explainability of learned policies, and predictive models, tested on physical robots.

\section{Research Threads}
\thread{Understanding People}{PRIOR WORK · 2021–2024}{How can a robot recognize what people are doing and what they mean from incomplete, multimodal evidence?}
\begin{compactitems}
\item Designed GCN–transformer models for action recognition and fall detection from multi-view human pose. Under occlusion, fusing views raised F1 from 0.76 to 0.82 (graduate course project). \evidence{https://tvpian.github.io/projects/multimodal-perception/}{PoseFusion}
\item Co-developed \href{https://arxiv.org/abs/2405.11655}{Track Anything Rapter}, in which a person picks a target by click, box, image, or text and a drone follows it using foundation-model perception and visual servoing.
\item Implemented an exploratory policy that learns from demonstration, speech, and gesture on the \href{https://tvpian.github.io/projects/natsgd/}{NatSGD dataset (Snehesh Shrestha, UMD)}. Whether gesture helps is not yet validated.
\end{compactitems}
\thread{Learning from People}{SYSTEM IMPLEMENTED; EVALUATION MANUAL}{How can robots improve from people's demonstrations and corrections without losing what they already do well?}
\begin{compactitems}
\item Built a workflow in which people demonstrate tasks, take over when the robot fails, and have their corrections reviewed and added back into training before a new policy is rolled out. \evidence{https://tvpian.github.io/\#continual-learning}{Overview}
\item Built a step-level control interface through which people and AI models operate the arm with the same small, safety-checked actions. Every step is recorded, so human and model runs can be compared and reused for training.
\item Developing a portable egocentric interface that records human demonstrations through the same data pipeline (interface in development). \evidence{https://tvpian.github.io/\#egocentric}{Overview}
\end{compactitems}
\thread{Policy Explainability and Representation}{EXPERIMENTS IN PROGRESS}{What has a policy learned, and what shapes its actions?}
\begin{compactitems}
\item Studied what drives learned policies using attention maps, Grad-CAM, input ablations, prompt comparisons, and self-supervised versus pretrained representations. These are diagnostics, not causal explanations. \evidence{https://tvpian.github.io/\#explainability}{Overview}
\end{compactitems}
\thread{Predictive Models, Inverse Dynamics, and Transfer}{OFFLINE EVALUATION COMPLETED}{How can predicted outcomes and inverse dynamics support action learning on physical robots?}
\begin{compactitems}
\item Adapted a published video-action architecture by retraining its action head on teleoperation data and evaluated it offline. Physical validation is next. \evidence{https://tvpian.github.io/\#world-models}{Overview}
\end{compactitems}
\section{Education}
\roleentry{Master of Engineering in Robotics, GPA: 3.9/4.0}{Aug. 2022 – May 2024}{University of Maryland, College Park}{Maryland, USA}
\par\vspace{3pt}\textbf{Selected coursework:} Advanced Techniques in Visual Learning and Recognition; Perception for Autonomous Robots; Planning for Autonomous Robots; Control of Robotic Systems; Fundamentals for Artificial Intelligence and Deep Learning Framework.
\par\vspace{4pt}
\roleentry{B.Tech. in Electronics and Communication Engineering, GPA: 8.59/10.0}{Jul. 2014 – Apr. 2018}{Vellore Institute of Technology}{India}
\newpage
\section{Research and Professional Experience}
\roleentry{Robotics Engineer – Robot Learning and Deformable Manipulation}{Jun. 2024 – Present}{General Motors R\&D}{Michigan, USA}
\begin{compactitems}
\item Architected and implemented a \href{https://tvpian.github.io/#current}{modular ROS 2 platform for bimanual robot learning}, from teleoperated demonstration and synchronized data capture to training, simulation, staged testing, and physical deployment. Shared interfaces let devices, learning methods, and simulators change without rebuilding the workflow.
\item Led the human-correction workflow, the step-level control interface, and the policy explainability and representation studies described under Research Threads.
\item Built foundation-model services for task-conditioned detection, segmentation, tracking, and pose servoing. Reusing embeddings and propagating masks over time cut perception latency by about 35\% and sustained 30–40 FPS.
\item Developed force-aware bimanual manipulation of deformable linear objects using impedance control and force-torque feedback. Without elastic-tension regulation, failures doubled.
\item Built a calibrated digital twin and a behavior-tree-driven synthetic-data workflow to test how policies degrade under changes in sensing, timing, contact, and appearance.
\item Designed the Cartesian motion-control and safety layer for learned policies, and supervised its extension to compliant (admittance-controlled) dual-arm data collection and rollout.
\item Compared action-chunking and open-source vision-language-action policies while developing evaluation measures beyond success rate (partial progress, failure type, recovery, repeatability).
\end{compactitems}
\roleentry{Researcher – AI and Robotics}{Feb. 2023 – Mar. 2024}{CATT Laboratory, University of Maryland}{Maryland, USA}
\begin{compactitems}
\item Architected the fusion pipeline for \href{https://doi.org/10.1007/978-3-031-74835-6_15}{multimodal anomaly detection}, with separate LiDAR, odometry, and network encoders and reconstruction-based scoring. The published system reached 98\% accuracy versus 72\% for LiDAR alone, and 97\% with half the training data.
\item Co-developed \href{https://arxiv.org/abs/2405.11655}{Track Anything Rapter}, a \href{https://github.com/tvpian/Project-TAR}{ROS 2 aerial system} for multimodal target specification (see Research Threads).
\end{compactitems}
\roleentry{Researcher – Collaborative Robotics and Computer Vision}{Jun. 2021 – Jan. 2023}{Indian Institute of Science and University of Maryland}{India and USA}
\begin{compactitems}
\item Developed a \href{https://www.researchgate.net/publication/399369200_Vision_Based_Hybrid_IK_Task_Planning_with_Feedforward_Neural_Network_for_Collaborative_Plant-Robot_Interaction_in_Precision_Farming}{vision-guided precision-weeding robot} that estimates plant roots and spread, then plans with hybrid inverse kinematics to reach weeds while reducing contact with crops.
\item Built synthetic-data and augmentation pipelines for weed localization. Hard-negative sampling and structured distractors improved robustness by 20–30\%.
\item Co-designed a \href{https://github.com/tvpian/Position_Based_Torque_Controller_for_Safe_Robotic_Weeding}{dual-purpose end effector} (plucking plus spraying or watering) and adaptive \href{https://drive.google.com/file/d/1IvKSo8U6aEWXKw8ucpV-9Hd1IqdPomFb/view?usp=sharing}{position/torque control} for delicate plants.
\item Implemented an exploratory \href{https://www.snehesh.com/natsgd/}{demonstration, speech, and gesture policy} (see Research Threads).
\item Developed and evaluated a \href{https://github.com/tvpian/Monocular_Depth_Estimation_using_Transfer_Learning}{U-Net monocular depth model} in an IROS 2021 challenge setting.
\end{compactitems}
\roleentry{Project Engineer – AI Innovation}{Jul. 2018 – May 2021}{CTO Office, Wipro Digital}{Bangalore, India}
\begin{compactitems}
\item Led applied-research prototypes in computer vision, HCI, recommendation, generative modeling, and mobile inference, and coordinated quarterly prototype selection and delivery.
\item Built tool tracking, mobile vehicle recognition, photogrammetry, \href{https://drive.google.com/file/d/1YryajvaO5w4Blw56dToFsm54rEQHvLkJ/view?usp=sharing}{human-pose estimation}, \href{https://github.com/tvpian/Virtual_Try-On}{virtual try-on}, and \href{https://github.com/tvpian/My-Style}{recommendation} components for industrial and retail workflows.
\end{compactitems}
\par\vspace{5pt}\textbf{Selected Earlier Technical Projects}\par
\projectentry{Implicit Neural Representation and SLIC Segmentation}{PyTorch, TensorFlow, OpenCV}{2023}
\begin{compactitems}
  \item Implemented \href{https://github.com/tvpian/Image-Outpainting}{coordinate-conditioned image outpainting} and \href{https://github.com/tvpian/SLIC-Image-Segmentation}{deep-feature superpixel segmentation}, including controlled comparisons with simpler baselines.
\end{compactitems}

\projectentry{Maze Runner and Autonomous Retrieval Robot}{ROS 2, Nav2, visual servoing}{2023}
\begin{compactitems}
  \item Built \href{https://github.com/tvpian/Maze_Runner}{visual-cue-based maze navigation} in \href{https://github.com/tvpian/Maze_Solver}{Gazebo} and an \href{https://youtu.be/5vNxQM7fV6A}{onboard-perception retrieval robot} using sensor fusion and visual servoing.
\end{compactitems}

\projectentry{Reinforcement-Learning Cricket Simulation}{adaptive agents and simulation}{2020}
\begin{compactitems}
  \item Built a simulated bowler agent that adapts to batting styles and field configurations. \evidence{https://github.com/tvpian/RL_Cricket_Simulation_Engine}{code}
\end{compactitems}

\newpage
\section{Publications}
\textit{(Also published as V. P. Tharun.)}

\begin{enumerate}[leftmargin=18pt,itemsep=2.4pt,topsep=2pt]
\item A. Zoulkarni, M. Noorani, J. S. Baras, J. Mirenzi, \textbf{T. V. Puthanveettil}, and C. D. Grody. ``\href{https://scholar.google.com/citations?view_op=view_citation\&hl=en\&user=Vre9wQQAAAAJ\&citation_for_view=Vre9wQQAAAAJ:aqlVkmm33-oC}{Efficient Integration of Domain Knowledge and Data-Driven Methods for Securing Autonomous Cyber-Physical Systems.}'' IEEE ROSE 2026. DOI: \href{https://doi.org/10.1109/ROSE69181.2026.11570974}{10.1109/ROSE69181.2026.11570974}.
\item \textbf{V. P. Tharun} and A. R. Chowdhury. ``\href{https://www.researchgate.net/publication/399369200_Vision_Based_Hybrid_IK_Task_Planning_with_Feedforward_Neural_Network_for_Collaborative_Plant-Robot_Interaction_in_Precision_Farming}{Vision-Based Hybrid IK Task Planning with Feedforward Neural Network for Collaborative Plant-Robot Interaction.}'' \textit{International Conference on Social Robotics + AI (ICSR + AI 2025)}, LNAI 16132, Springer. DOI: \href{https://doi.org/10.1007/978-981-95-2382-5_18}{10.1007/978-981-95-2382-5\_18}.
\item M. Noorani, \textbf{T. V. Puthanveettil}, A. Zoulkarni, J. Mirenzi, C. D. Grody, and J. S. Baras. ``\href{https://doi.org/10.1007/978-3-031-74835-6_15}{Multimodal Anomaly Detection for Autonomous Cyber-Physical Systems Empowering Real-World Evaluation.}'' \textit{GameSec}, 2024.
\item L. P. Pratap, P. M. Shailendrasingh, A. Anand, and \textbf{V. P. Tharun}. ``\href{https://ieeexplore.ieee.org/abstract/document/8392038}{Wall-Climbing Robot Using Soft Robotics.}'' \textit{ICPCSI}, 2017.
\end{enumerate}
Earlier work in data mining and forecasting (2018–2020): five papers and reports on rainfall forecasting and water-quality assessment. \evidence{https://scholar.google.com/citations?user=Vre9wQQAAAAJ\&hl=en}{Google Scholar}
\section{Preprints and Research Outputs}

\begin{enumerate}[leftmargin=18pt,itemsep=2.2pt,topsep=2pt]
\item \textbf{T. V. Puthanveettil}, A. Singh, Y. Jain, V. Bukka, and S. Arjun S. ``\href{https://arxiv.org/abs/2405.11659}{Auto-Platoon: Freight by Example.}'' \textit{arXiv:2405.11659}, 2024.  \evidence{https://github.com/tvpian/Auto-Platoon}{code}
\item \textbf{T. V. Puthanveettil} and F. Obaid ur Rahman. ``\href{https://arxiv.org/abs/2405.11655}{Track Anything Rapter (TAR).}'' \textit{arXiv:2405.11655}, 2024.  \evidence{https://github.com/tvpian/Project-TAR}{code + demos}
\item \textbf{Tharun Puthanveettil}, Abhijay Singh, Aniruddh Balram. ``\href{https://www.researchgate.net/publication/378143241_Pose_Fusion_-_Multi-View_Pose_Integration_for_Comprehensive_Action_Recognition}{Pose Fusion - Multi-View Pose Integration for Comprehensive Action Recognition.}'' UMD graduate course project. Project poster and preprint.  \evidence{https://github.com/tvpian/HPE_for_HAR/blob/main/preprint.pdf}{preprint}
\item \textbf{V. P. Tharun} et al. ``\href{https://drive.google.com/file/d/1piWRkBv1j-_5-h8RNXS63Tf8b7eT6lat/view?usp=sharing}{Application of Mobile Collaborative Robot using Deep Learning in Precision Weed Control.}'' Manuscript, 2021.
\item \textbf{V. P. Tharun} et al. ``\href{https://www.researchgate.net/publication/378143055_Event_Based_Dynamic_Obstacle_Avoidance_For_Outdoor_Environments}{Event-Based Dynamic Obstacle Avoidance in Outdoor Environments.}'' IROS 2021 workshop poster.
\end{enumerate}

\section{Technical Skills}
\textbf{Robot learning:} PyTorch, TensorFlow, LeRobot, action-chunking policies, vision-language-action models, diffusion/flow policies, imitation learning, self-supervised learning\par\vspace{2pt}
\textbf{Robotics and control:} ROS 2, ROS, MoveIt 2, impedance/compliance control, visual servoing, behavior trees, bimanual manipulation, force-torque feedback\par\vspace{2pt}
\textbf{Simulation and perception:} Isaac Sim/Lab, MuJoCo, Gazebo, OpenCV, Open3D, RGB-D cameras, detection, segmentation, tracking, pose estimation\par\vspace{2pt}
\textbf{Programming and systems:} Python, C++, SQL, Linux, Docker, Git, REST APIs
\section{Intellectual Property, Mentoring, and Communication}
\textbf{Patents:} Inventor on three filed patent applications; details withheld pending publication.\par\vspace{3pt}
\textbf{Research supervision:} Supervised an engineering intern extending the robot-learning platform for compliant bimanual data collection, learning, analysis, and rollout.\par\vspace{3pt}
\textbf{FIRST Robotics:} Electrical Programming Head Coach, The Irrational Engineers (Jun. 2024–Present); mentor high-school students in controls, computer vision, and software.\par\vspace{3pt}
\textbf{Writing:} \href{https://tvpian.github.io/human-in-the-loom/}{\textit{Human in the Loom}}, essays on robot learning, embodied intelligence, and credible research practice. \evidence{https://tvpian.substack.com/archive}{Essay archive}
\end{document}
